\documentclass[12pt]{article}
\usepackage[utf8]{inputenc}
\usepackage[T1]{fontenc}
\usepackage{hyperref}
\title{Sukuma Language Resources for K-12 Curriculum Development}
\author{Compiled by Claude Sonnet 5.5 (claude-sonnet-5-5) for Andr\'as Kornai}
\date{}
\begin{document}
\maketitle

\section{Introduction}
Sukuma (Kisukuma, ISO 639-3 \texttt{suk}) is a Bantu language of Tanzania, spoken southeast of Lake Victoria between Mwanza, Shinyanga and Lake Eyasi. Wikipedia gives a figure of about 8.1 million speakers, higher than the 5.4 million in the survey description. The main dialects are generally mutually intelligible. It uses a Roman orthography similar to Swahili's.

\section{Primary Educational Resources}
No Sukuma primers, textbooks or school programmes were found in the searches.

\section{Linguistic and Reference Materials}
\begin{itemize}
\item Wikipedia, ``Sukuma language''. \url{https://en.wikipedia.org/wiki/Sukuma_language}
\item Omniglot, ``Sukuma (Kisukuma)''. It describes the writing system and includes a sample of the Lord's Prayer. \url{https://omniglot.com/writing/sukuma.htm}
\item Historical African Languages Database (BYU). It catalogues A. Seidel's 1894 German--Kisukuma dictionary and grammatical sketch. \url{https://hald.byu.edu/summary_sukuma_Seidel}
\item REFLEX (CNRS) bibliography for Sukuma. It lists many dictionaries and grammars, including a Kisukuma--English dictionary, an English--Kisukuma dictionary, a French--Sukuma dictionary, ``Grammar of the Kisukuma language'', K\"onen's ``New KiSukuma grammar'', Downey's ``KiSukuma grammar'' and Augustiny (1929). These are catalogue entries only, and I did not verify that any can be obtained. \url{https://reflex.cnrs.fr/Lexiques/webball/biblio.php?LGNOM=Suk}
\end{itemize}

\section{Religious and Community Resources}
\begin{itemize}
\item Scripture Earth, Sukuma. It lists Bible text and audio, including a 1997 Bible League edition and a Common Language Version. \url{https://www.scriptureearth.org/00eng.php?iso=suk}
\end{itemize}

\section{Digital and Audio Resources}
Scripture Earth (above) links to Bible audio and text. No other audio, corpus or software resources were found.

\section{Curriculum Development Context}
Sukuma has a large speaker base, a Roman orthography, full or partial scripture and a considerable older linguistic literature. It is not known to be taught in schools, and I found no school materials. A curriculum would need to build reading materials from the existing orthography and scripture texts. The older dictionaries and grammars would need to be located and checked for currency.
\end{document}